\section{Introduction}
\label{sec:research_method_introduction}

This chapter outlines the methodological framework used to investigate the
research goals and hypotheses introduced earlier. Prior work on compositional
generative modelling—whether score-based, energy-based, or density-based—offers
mechanisms for combining pretrained diffusion models but lacks a principled
understanding of when such operations yield semantically meaningful
compositions. Score arithmetic often fails under entangled or misaligned
factors, Product-of-Experts methods rely on factorisation assumptions that seldom
hold, and density-based approaches such as SuperDiff provide mathematically
correct density combinations without guaranteeing semantic correctness.

Our methodology addresses these gaps through a three-phase structure. Phase~1
evaluates how existing compositional operators behave in a controlled medical
testbed, identifying when they succeed or fail to produce plausible hybrid
pathologies. Phase~2 extends composition to property-guided settings, testing
whether continuous attributes such as severity or localisation can be reliably
controlled. Phase~3 employs Riemannian diffusion analysis to examine the latent
and geometric conditions that underlie compositional correctness and explain the
failure modes observed earlier.

Together, these phases provide a unified empirical and structural framework for
understanding and improving semantic compositionality in generative modelling.